\section{Discussion and Conclusion}
\label{sec:conclusion}

In this work, we presented a rigorous methodological dissection of Sharpness-Aware Isotropic Merging (SAIM), a recent dual-stage framework that combines a custom coordinate-wise optimizer (SA-BCD) with SVD-based adaptive isotropic merging. By cross-evaluating five optimizers and three merging strategies on a comprehensive $5 \times 3$ grid, we systematically isolated the true causal drivers of performance in continual model merging.

Our empirical findings offer a sobering lesson on the inflation of complex frameworks in deep learning. We demonstrated that:
\begin{enumerate}
  \item \textbf{The benefit of post-hoc SVD merging is highly sensitive to boundary conditions.} Rather than being a universal remedy for representation collapse, we demonstrate that SVD-based isotropic merging's regularizing effect only manifests under active parameter-mixing regimes (such as $\lambda = 0.2$), where it mitigates parameter interference. Under standard sequential fine-tuning parity ($\lambda = 0.0$), applying SVD interpolation to the expert update is redundant and actively distorts un-mixed parameters, degrading average accuracy by $3\%$ to $12\%$ across all optimizers tested.
  \item \textbf{Optimizer-driven flatness is the foundational driver of merging performance.} Fine-tuning task-specific experts using standard, globally-perturbed Sharpness-Aware Minimization (SAM) allows their weights to be consolidated without encountering high-loss barriers. This simple baseline (\emph{SAM} + \emph{Task Arithmetic}) yields an impressive \textbf{68.27\%} average accuracy under $\lambda=0.0$, and provides the critical foundation for SVD's success under $\lambda=0.2$.
  \item \textbf{The proposed SA-BCD optimizer is mathematically flawed and suboptimal.} We identified a fatal algebraic bug in SAIM's published formula that causes optimization divergence. Furthermore, restricting perturbations to a selected coordinate subset is empirically suboptimal compared to global perturbations.
\end{enumerate}

These results lead to a clear, actionable recommendation for practitioners: rather than employing expensive, multi-component post-hoc weight manipulations on un-mixed models, practitioners should focus on \textbf{optimizing for flatness during individual model training} (e.g., using standard SAM). When active weight consolidation is required under non-zero mixing regimes, post-hoc weight consolidation (such as SVD isotropic merging) can be highly beneficial, but it must be built on top of a flat parameter foundation to realize its full potential.

More broadly, our study underscores a critical message for the machine learning community: \emph{rigorous ablation is mandatory.} We strongly urge authors of future model merging works to perform sequential, component-by-component evaluations, compare their methods against heavily-tuned standard baselines (such as standard SAM), and ensure that mathematical formulas match their underlying code implementations. Only through such disciplined methodological audits can we prevent false progress and establish a solid, transparent foundation for future research in deep model merging. To support future model-merging audits and reproducibility, we will release our complete modular $5 \times 3$ grid evaluation suite and our corrected, verified implementations of the SA-BCD optimizer as public, open-source artifacts.
